% !TEX spellcheck = en_US
% !TEX spellcheck = LaTeX
\documentclass[a4paper,english,10pt]{article}
\usepackage{%
amsfonts,%
amsmath,%
amssymb,%
amsthm,%
babel,%
bbm,%
%biblatex,%
caption,%
centernot,%
color,%
enumerate,%
%enumitem,%
epsfig,%
epstopdf,%
etex,%
fancybox,%
framed,%
fullpage,%
%geometry,%
graphicx,%
hyperref,%
latexsym,%
mathptmx,%
mathtools,%
multicol,%
paralist,%
pgf,%
pgfplots,%
pgfplotstable,%
pgfpages,%
proof,%
psfrag,%
%subfigure,%
tcolorbox,%
tfrupee,%
tikz,%
times,%
%ulem,%
url,%
xcolor,%
mathpazo,%
}

\definecolor{shadecolor}{gray}{.95}%{rgb}{1,0,0}
\usepackage[margin=1in,top=0.75in]{geometry}
\usepackage[mathscr]{eucal}
\usepgflibrary{shapes}
\usepgfplotslibrary{fillbetween}
\usetikzlibrary{%
arrows,%
backgrounds,%
chains,%
decorations.pathmorphing,% /pgf/decoration/random steps | erste Graphik
decorations.text,% 
matrix,%
positioning,% wg. " of "
fit,%
patterns,%
petri,%
plotmarks,%
scopes,%
shadows,%
shapes.misc,% wg. rounded rectangle
shapes.arrows,%
shapes.callouts,%
shapes%
}

%\pgfplotsset{compat=newest} %<------ Here
\pgfplotsset{compat=1.11} %<------ Or use this one

\theoremstyle{plain}
\newtheorem{thm}{Theorem}[section]
\newtheorem{lem}[thm]{Lemma}
\newtheorem{prop}[thm]{Proposition}
\newtheorem{cor}[thm]{Corollary}
\newtheorem{clm}[thm]{Claim}

\theoremstyle{definition}
\newtheorem{defn}[thm]{Definition}
\newtheorem{conj}[thm]{Conjecture}
\newtheorem{exmp}[thm]{Example}
\newtheorem{axiom}[thm]{Axiom}
\newtheorem{assum}[thm]{Assumption}
\newtheorem{exerc}[thm]{Exercise}
\newtheorem*{defin}{Definition}

\theoremstyle{remark}
\newtheorem{rem}{Remark}
\newtheorem{note}{Note}

\newcommand{\iid}{\emph{i.i.d.}\ }
\newcommand{\Cov}{\operatorname{Cov}}
%\newcommand{\det}{\operatorname{det}}
%\newcommand{\Real}{\mathbb{R}}
\newcommand{\tr}{\operatorname{tr}}
\newcommand{\Var}{\operatorname{Var}}
\newcommand{\cov}{\operatorname{cov}}

%\renewcommand{\proof}[1]{\begin{proof}#1\end{proof}}
\newcommand{\EQ}[1]{\begin{equation*}#1\end{equation*}}
\newcommand{\EQN}[1]{\begin{equation}#1\end{equation}}
\newcommand{\eq}[1]{\begin{align*}#1\end{align*}}
\newcommand{\meq}[2]{\begin{xalignat*}{#1}#2\end{xalignat*}}
\newcommand{\norm}[1]{\left\lVert#1\right\rVert}
\newcommand{\inner}[1]{\left\langle #1\right\rangle}
\newcommand{\abs}[1]{\left\lvert#1\right\rvert}
\newcommand{\expect}[1]{\mathbb{E}\left[{#1}\right]}
\newcommand{\prob}[1]{\mathbb{P}\left[{#1}\right]}
\newcommand{\given}{\; \big\vert \;} 
\newcommand{\set}[1]{\left\{#1\right\}} 
\newcommand{\indicator}[1]{1_{\set{#1}}} 
\newcommand{\Ind}[1]{\mathbbm{1}_{#1}} 
\newcommand{\SetIn}[1]{\mathbbm{1}_{\set{#1}}} 
\newcommand{\red}[1]{\textcolor{red}{#1}} 
\newcommand{\floor}[1]{\left\lfloor#1\right\rfloor}
\newcommand{\ceil}[1]{\left\lceil#1\right\rceil}


\newcommand{\C}{\mathbb{C}}
\newcommand{\D}{\mathbb{D}}
\newcommand{\E}{\mathbb{E}}
\newcommand{\F}{\mathbb{F}}
\newcommand{\N}{\mathbb{N}}
\renewcommand{\P}{\mathbb{P}}
\newcommand{\Q}{\mathbb{Q}}
\newcommand{\R}{\mathbb{R}}
\newcommand{\Z}{\mathbb{Z}}

\newcommand{\bA}{\mathbf{A}}
\newcommand{\bI}{\mathbf{I}}
\newcommand{\bU}{\mathbf{1}}
\newcommand{\bx}{\mathbf{x}}
\newcommand{\bX}{\mathbf{X}}
\newcommand{\bZ}{\mathbf{Z}}


\newcommand{\cA}{\mathcal{A}}
\newcommand{\cB}{\mathcal{B}}
\newcommand{\cC}{\mathcal{C}}
\newcommand{\cF}{\mathcal{F}}
\newcommand{\cG}{\mathcal{G}}
\newcommand{\cM}{\mathcal{M}}
\newcommand{\cN}{\mathcal{N}}
\newcommand{\cP}{\mathcal{P}}
\newcommand{\cT}{\mathcal{T}}
\newcommand{\cX}{\mathcal{X}}
\newcommand{\cY}{\mathcal{Y}}

\newcommand{\sA}{\mathscr{A}}
\newcommand{\sB}{\mathscr{B}}
\newcommand{\sC}{\mathscr{C}}
\newcommand{\sE}{\mathscr{E}}
\newcommand{\sF}{\mathscr{F}}
\newcommand{\sG}{\mathscr{G}}
\newcommand{\sH}{\mathscr{H}}
\newcommand{\sL}{\mathscr{L}}
\newcommand{\sS}{\mathscr{S}}
\newcommand{\sT}{\mathscr{T}}
\newcommand{\sX}{\mathscr{X}}
\newcommand{\sY}{\mathscr{Y}}
\newcommand{\sZ}{\mathscr{Z}}

\newcommand{\tS}{\tilde{S}}
% Debug
\newcommand{\todo}[1]{\begin{color}{blue}{{\bf~[TODO:~#1]}}\end{color}}

% a few handy macros

\renewcommand{\le}{\leqslant}
\renewcommand{\ge}{\geqslant}
\newcommand\matlab{{\sc matlab}}
\newcommand{\goto}{\rightarrow}
\newcommand{\bigo}{{\mathcal O}}
%\newcommand{\half}{\frac{1}{2}}
%\newcommand\implies{\quad\Longrightarrow\quad}
\newcommand\reals{{{\rm l} \kern -.15em {\rm R} }}
\newcommand\complex{{\raisebox{.043ex}{\rule{0.07em}{1.56ex}} \hskip -.35em {\rm C}}}


% macros for matrices/vectors:

% matrix environment for vectors or matrices where elements are centered
\newenvironment{mat}{\left[\begin{array}{ccccccccccccccc}}{\end{array}\right]}
\newcommand\bcm{\begin{mat}}
\newcommand\ecm{\end{mat}}

% matrix environment for vectors or matrices where elements are right justifvied
\newenvironment{rmat}{\left[\begin{array}{rrrrrrrrrrrrr}}{\end{array}\right]}
\newcommand\brm{\begin{rmat}}
\newcommand\erm{\end{rmat}}

% for left brace and a set of choices
%\newenvironment{choices}{\left\{ \begin{array}{ll}}{\end{array}\right.}
\newcommand\when{&\text{if~}}
\newcommand\otherwise{&\text{otherwise}}
% sample usage:
%\delta_{ij} = \begin{choices} 1 \when i=j, \\ 0 \otherwise \end{choices}


% for labeling and referencing equations:
\newcommand{\eql}{\begin{equation}\label}
\newcommand{\eqn}[1]{(\ref{#1})}
% can then do
%\eql{eqnlabel}
%...
%\end{equation}
% and refer to it as equation \eqn{eqnlabel}.
% some useful macros for finite difference methods:
\newcommand\unp{U^{n+1}}
\newcommand\unm{U^{n-1}}
% for chemical reactions:
\newcommand{\react}[1]{\stackrel{K_{#1}}{\rightarrow}}
\newcommand{\reactb}[2]{\stackrel{K_{#1}}{~\stackrel{\rightleftharpoons}
 {\scriptstyle K_{#2}}}~}
\makeatletter
\def\th@plain{%
\thm@notefont{}% same as heading font
\itshape % body font
}
\def\th@definition{%
\thm@notefont{}% same as heading font
\normalfont % body font
}
\makeatother
\date{}



%opening
\title{Lecture-28: Compound Poisson Processes}
\author{}
\date{}
\begin{document}
\maketitle

\section{Compound Poisson process}
\begin{defn}
A \textbf{compound Poisson process} is a real-valued right-continuous process $Z:\Omega\to\R_+^{\R_+}$ with the following properties.
\begin{enumerate}[i\_]
	\item \textbf{Finitely many jumps:} for all $\omega \in \Omega$, sample path $t \mapsto Z_t(\omega)$ has finitely many jumps in finite intervals,
	\item \textbf{Independent increments:} for all $t,s \ge 0$; increments $Z_{t+s}-Z_t$ is independent of past $\sF_t \triangleq \sigma(Z_u: u\le t)$, 
	\item \textbf{Stationary increments:} for all $t,s \ge 0$, distribution of $Z_{t+s}-Z_t$ depends only on $s$ and not on $t$.
\end{enumerate}	
\end{defn}
\begin {defn}
For each $\omega \in \Omega$ and $n \in \N$, we can define time and size of $n$th jump
\meq{2}{
&\tS_0(\omega) = 0,&&\tS_n(\omega) = \inf\set{t > \tS_{n-1} : Z_t(\omega) \neq Z_{\tS_{n-1}}(\omega)}\\
&Y_0(\omega) = 0, &&Y_n(\omega) = Z_{\tS_n}(\omega) - Z_{\tS_{n-1}}(\omega).
}
\end{defn}
\begin{rem}
Recall that $\sF_s = \sigma(Z_u, u \in (0, s])$ is the collection of historical events until time $s$ 
%, and $\sF_{\bullet} = (\sF_s: s \ge 0)$ be the natural filtration 
associated with the process $Z$. 
If $\tS_n$ is almost surely finite for all $n \in \N$, 
then the sequence of jump times $\tS:\Omega\to\R_+^\N$ is a sequence of stopping times with respect to the natural filtration $\sF_\bullet$ of the process $Z$. 
\end{rem}
\begin{rem}
Let $N:\Omega\to\Z_+^{\R_+}$ be the simple counting process associated with the number of jumps of compound Poisson process $Z$ in $(0,t]$ defined by $N_t \triangleq \sum_{n\in\N}\SetIn{\tS_n\le t}$ for all $t\in\R_+$. 
Then, $\tS_n$ and  $Y_n$ are the respectively the arrival instant and the size of the $n$th jump, 
and we can write $Z_t = \sum_{i = 1}^{N_t}Y_i$. 
\end{rem}
%From finite jump property, it follows that $\set{S_n < \infty}$ almost surely for each $n \in \N$. 

%filtration $\sF_{\bullet}$. 
%\begin{defn}
\begin{prop}
A stochastic process $Z:\Omega\to\R_+^{\R_+}$ is a compound Poisson process iff its jump times form a Poisson process and the jump sizes form an \iid random sequence independent of the jump times. 
%can be represented as $Z_t=\sum_{i=1}^{N_t}X_i$ for all $t\geq 0$ where $\set{N_t,t\geq 0}$ is Poisson process and $\set{X_i, i\in \N}$ are \iid random variables independent of $\set{N_t, t\geq 0}$.
\end{prop}
\begin{proof} 
We will prove it in two steps. 
\begin{itemize} 
\item[Implication:] 
Let $Z$ be a compound Poisson process with the jump instant sequence $\tS:\Omega\to\R_+^\N$ and the jump size sequence $Y:\Omega\to\R_+^\N$. 
We will show that the counting process $N:\Omega\to\Z_+^{\R_+}$ is simple and has stationary and independent increments and the jump size sequence $Y$ is \iid. 
\begin{itemize}
%\item[Counting:] It is clear that the simple counting process $N$ can be completely determined by $\sF_t = \sigma(Z_u: u \le t)$. 
%, i.e. $N_t \in \sF_t$.  
\item[Independence of jumps and increments:] 
From the definition of jump instant sequence $\tS$, it follows that the counting process $N$ is adapted to the natural filtration $\sF_\bullet$ of the compound Poisson process $Z$. 
Since $Z_{t+s} - Z_t = \sum_{i = N_t+1}^{N_{t+s}}Y_i$, and the compound Poisson processes have independent increments, 
it follows that the increment $(N_{t+s}-N_t: s \ge 0)$ and $(Y_{N_t+j}: j \in \N)$ are independent of the past $\sF_t$. %$(Z_{u}: u \le t)$. 
\item[Stationarity:] Let's assume that step sizes are positive, then
we have 
\meq{2}{
&\tS_n = \inf\set{t > \tS_{n-1}: Z_t > Z_{\tS_{n-1}}}, &\text{ and }&\set{N_{t+s} - N_t = 0} = \set{Z_{t+s} - Z_t = 0}.
}  
From the stationarity of the increments it follows that the probability $P\set{N_{t+s}-N_t = 0}$ is independent of $t$ and equal to $e^{-\lambda s}$ for some $\lambda \in \R_+$.  
It follows that the counting process $N:\Omega\to\Z_+^{\R_+}$ has stationary increments and the associated jump sequence $\tS$ is homogeneous Poisson with intensity density $\lambda$.  

\item[Strong Markovity:] The compound Poisson process has the Markov property from stationary and independent increment property. 
Further, since each sample path $t \mapsto Z_t$ is right continuous, the process satisfies the strong Markov property at each almost sure stopping time. 

\item[Inter jump times \iid:]
We will inductively show that $\tS:\Omega\to\R_+^\N$ is a stopping time sequence and hence the inter jump times sequence $X:\Omega\to\R_+^\N$ defined by $X_n \triangleq \tS_n-\tS_{n-1}$ for each $n\in\N$ is an \iid sequence. 
From the exponential distribution of $\tS_1$, it follows that it is almost surely finite and hence a stopping time. 
From the stationary of increments of compound Poisson process and the strong Markov property at stopping times $\tS_2 -\tS_1 = \tS_1$ is independent of $\sF_{\tS_1}$ and identical in distribution to $\tS_1$. 
The result follows inductively. 

%Further,  since $S_n$ are almost sure stopping times with respect to the compound Poisson process, 
%it also has strong Markov property from the right continuity of the sample paths. 
\item[Jump size \iid: ] 
From strong Markov property of $Z$, the jump size $Y_n$ is independent of the past $\sF_{\tS_{n-1}} = \sigma(Z_u: u \le \tS_{n-1})$ and from stationarity it is identically distributed to $Y_1$ for each $n \in \N$. 
It follows that the jump size sequence $Y$ is \iid and independent of jump instant sequence $\tS$.  
\item[Superposition:] Similar arguments can be used to show for negative jump sizes. 
For real jump sizes, we can form two independent Poisson processes with negative and positive jumps, 
and the superposition of these two processes is Poisson. 
\end{itemize}

\item[Converse:] Let $X:\Omega\to\R_+^\N$ be an \iid inter-jump sequence distributed exponentially with rate $\lambda$ and $Y:\Omega\to\R^\N$ be an \iid jump size sequence independent of $X$. 
We can define the jump instant sequence $\tS:\Omega\to\R_+^\N$ defined as $\tS_n \triangleq \sum_{i=1}^nX_i$ for each $n\in\N$, 
the counting process for the number of jumps $N:\Omega\to\Z_+^{\R_+}$ defined as $N_t \triangleq \sum_{n\in\N}\SetIn{\tS_n \le t}$ for each $t\in\R_+$, 
and the compound process $Z:\Omega\to\R^{\R_+}$  defined as $Z_t \triangleq \sum_{n =1}^{N_t}Y_n$ for each $t\in\R_+$. 
\begin{itemize}
\item[Finitely many jumps:] Since $N_t$ is finite for any finite $t$, it follows that the compound Poisson process $Z$ has finitely many jumps in finite intervals. 
\item[Independence of increments:] For any finite $n \in \N$ and finite intervals $I_i$ for $i \in [n]$, 
we can write $Z(I_i)= \sum_{k=1}^{N(I_i)}Y_{ik}$, where $Y_{ik}$ denotes the $k$th jump size in the interval $I_i$. 
Since the independent sequence $(N(I_i): i \in [n])$ and $Y:\Omega\to\R_+^\N$ are also mutually independent, 
it follows that $Z(I_i)$ are independent. 
\item[Stationarity of increments:] Further, the stationarity of the increments of the compound process is inferred from the distribution of $Z(I_i)$, which is 
\EQ{
P\set{Z(I_i) \le x} 
= \sum_{m \in \Z_+}P\set{Z(I_i) \le x, N(I_i) = m}  
= \sum_{m \in \Z_+}P\set{\sum_{k=1}^mY_{ik} \le x}P\set{N(I_i) = m}. 
}
\end{itemize}
\end{itemize}
\end{proof}
%\end{defn}
%Alternately it can also be defined in the following way.
%%\begin{defn}
%A compound Poisson Process is a point Process \set{$Z_t,t \geq 0$} having the following properties.
%\begin{enumerate}
%	\item For all $\omega \in \Omega, t \longmapsto Z_t(\omega)$ has finitely many jumps in finite intervals.
%	\item For all $t,s \geq 0; Z_{t+s}-Z_t$ is independent of $\set{Z_u, u\leq t}$.
%	\item For all $t,s \geq 0$, distribution of $Z_{t+s}-Z_t$ depends only on $s$ and not on $t$.
%\end{enumerate}	
%%\end{defn}
%\begin{defn} 
%A compound Poisson Process is stationary and independent increments point Process with jump points $S_n=\inf \set{t>0 | N(t)=n}$, and associated jump sizes $X_n$ independent of jump instants.% are the jump sizes associated with the Process.
%\end{defn}

\begin{shaded*}
\begin{exmp}  Examples of compound Poisson processes. 
\begin{itemize}
\item Arrival of customers in a store is a Poison process $N$. 
Each customer $i$ spends an \iid amount $X_i$ independent of the arrival process. 
\meq{2}{
&Y_0=0,&&Y_n=\sum_{i=1}^{n}X_i, i \in [n].
}
Now define $Z_t \triangleq Y_{N_t}$ as the amount spent by the customers arriving until time $t \in \R_+$. 
Then $Z:\Omega\to\R_+^{\R_+}$ is a compound Poisson Process.
%\end{exmp}
%\begin{exmp}  
\item Let the time between successive failures of a machine be independent and exponentially distributed. 
The cost of repair is \iid random at each failure. 
Then the total cost of repair in a certain time $t$ is a compound Poisson Process. 
\end{itemize}
\end{exmp}
\end{shaded*}

\end{document}
\section{Generalizations}
\begin{defn}
A compound Poisson point process $Z:\Omega\to\R^{\sB(\sX)}$ is defined for a Poisson point process $S: \Omega\to\sX^\N$ and an \iid process $Y:\Omega\to\R^\sX$, for each bounded $A \in \sB(\sX)$ as 
\EQ{
Z(A) \triangleq \sum_{n \in \N}Y(S_n)\Ind{A}(S_n).
}
\end{defn}
\begin{thm} 
For any finite $k \in \N$, a finite sequence of bounded disjoint sets $A \in \sB(\sX)^k$, and a vector $z \in \R^k$, the finite dimensional distribution of the random vector $(Z(A_1), \dots, Z(A_k))$ is given in terms of $B\triangleq \cup_{i=1}^kA_i$, as
\EQ{ 
P(\cap_{i=1}^k\set{Z(A_i) \le z_i}) 
= \prod_{i=1}^kP\set{Z(A_i) \le z_i} 
= e^{-\Lambda(B)}\prod_{i=1}^k\sum_{n_i \in \Z_+}\frac{\Lambda(A_i)^{n_i}}{n_i!}P\set{\sum_{j=1}^{n_i}Y_{i,j} \le z_i}.
}
\end{thm}
\begin{thm} 
For any Borel measurable function $f:\sX\to\R_+$, we can write the Laplace functional of compound Poisson process $Z$ as 
\EQ{
\sL_Z(f) = \exp\left(-\int_{\sX}(1- \E e^{-f(x)Y(x)})d\Lambda(x)\right).
}
\end{thm}
\begin{proof}
Let $A \in \sB(\sX)$ be a bounded set, then we can write the Laplace functional for function $g \triangleq f\Ind{A}$ as 
\EQ{
\sL_Z(g) 
= \expect{\expect{\exp\Big(-\int_Af(x)dZ(x)\Big)\mid N(A)}}.
}
Using the fact that points $S\cap A$ are \iid with density $\frac{\lambda(x)}{\Lambda(A)}$ conditioned on $\set{N(A) = n}$, and $Y$ is independent, we get 
\EQ{
\expect{\exp\Big(-\int_Af(x)dZ(x)\Big)\mid \set{N(A)=n}} 
= \expect{\prod_{x \in S\cap A}e^{-f(x)Y(x)}\mid \set{N(A)=n}}
= \Big(\int_{x \in A} \expect{e^{-f(x)Y(x)}}\frac{\lambda(x)}{\Lambda(A)}\Big)^n.
}
Since $N(A)$ is a discrete random variable with a Poisson probability mass function with rate $\Lambda(A)$, we get the result by taking limits of increasing bounded sets $A_k\uparrow \sX$ and monotone convergence theorem. 
\end{proof}
\end{document}

\section{Non-stationary Poisson process}
From the characterization of Poisson process just stated, we can generalize to non-homogeneous Poisson Process. 
In this case, the rate of Poisson Process $\lambda$ is time varying. 
%It is not clear from the first two characterizations, how to generalize the definition of Poisson process to the non-homogeneous case. 
%We used third characterization of Poisson process for this generalization. 
A simple counting process $(N(t): t\ge 0)$ is said to be possibly \textbf{non-stationary Poisson process} if $N(0) =0$ and it has independent increments. 
That is,
\begin{enumerate}[i\_]
\item \textbf{simple counting process:} for each $\omega \in \Omega$, each sample path $t \mapsto N_t(\omega)$ is right continuous, non-decreasing, integer valued, $N(0) = 0$, and has jumps of unit size only,
\item \textbf{independent increments:} for any $t,s \ge 0$, the random variable $N_{t+s} - N_t$ is independent of the past $(N_u: u \le t)$. 
\end{enumerate}
	
Let $\Lambda(t) = \E N_t$ for all $t \ge 0$. 
From non-decreasing property of counting processes, it follows that the mean is also non-decreasing in time $t$. 
From right continuity of counting process and the monotone convergence theorem, it follows that mean function is also right continuous. 
The \textbf{time inverse} of mean is defined as
\EQ{
\tau(t) = \inf\set{ s > 0 : \Lambda(s) > t},~ t \ge 0. 
}
It follows that the following events are identical, $\set{s > \tau(t)} = \set{\Lambda(s) > t}$.  
%Further, for all $s < \tau(t)$, we have $\Lambda(s) \le t$. 
Therefore, if $\Lambda$ is a continuous function, we would have $\Lambda(\tau(t)) = t$. 
Since inverse of a non-decreasing function is also non-decreasing, 
we conclude that $\tau(t)$ is non-decreasing function of time $t$. 
We can also see that by taking $t_1 \le t_2$. 
Then for all $s > \tau(t_2)$ we have $\Lambda(s) > t_2 \ge t_1$. 
In particular, we have $\Lambda(s) > \tau(t_1)$ and hence $\tau(t_2) \ge \tau(t_1)$. 
From the definition of $\tau(t)$, it follows that $\set{s > \tau(\Lambda(t))}= \set{\Lambda(s) > \Lambda(t)}$. 
Hence, 
\EQ{
\tau(\Lambda(t)) = \sup\set{s > 0: \Lambda(s) \le \Lambda(t)}. 
}

If $\Lambda$ is differentiable for all $t > 0$, then we can define $\lambda(t) = \Lambda'(t)$ for all $t > 0$. 
In higher dimensions, we have $\Lambda(A) = \E N(A)$ for some Borel measurable set $A \in \sB(\R^d)$. 
If $\Lambda$ is differentiable for all $x \in \R^d$, then for $A = \prod_{i=1}^d(x_i, x_i+h_i]$ we can write the intensity density $\lambda(x) = \lim_{h_1, \dots, h_d \to 0}\frac{m(A)}{h_1\dots h_d}$. 
In other words, we have $\Lambda(A) = \E N(A) = \int_{x \in A}\lambda(x)dx$. 

\begin{thm} 
Let $N_t$ be a non-stationary Poisson process, such that $\Lambda(t) = \E N_t$ is continuous. 
Then, the process $M: \Omega \times \R_+ \to \Z_+$ is a stationary Poisson process with unit rate, 
defined point-wise as 
\EQ{
M_t(\omega) \triangleq N_{\tau(t)}(\omega),~~t \ge 0, \omega \in \Omega. 
}
\end{thm}
\begin{proof} 
Since $\tau$ is non-decreasing and continuous and $N$ is a simple counting process, 
it follows that $M$ is right continuous, non-decreasing, integer valued, and simple.    
Fix $t > s \ge 0$ and let $s' \triangleq \tau(s)$ and $t' \triangleq \tau(t)$.  
Then, by definition of $M_t, t', s'$ and independent increment property of non-stationary Poisson process $N$, 
we have $M_t - M_s = N_{t'}-N_{s'}$ is independent of the past $(N_u: u \le s') = (M_u: u \le s)$. 
That is, the process $M$ has independent increments. 
For stationarity, we see that mean rate of increment is unity and stationary, 
%\EQ{
%P\set{M_t - M_s = 0} = 
%}
\EQ{
\E[M_{t}-M_s | M_u; u \le s] = \E[N_{t'}-N_{s'} | N_u; u \le s'] = \Lambda(t') - \Lambda(s') = \Lambda(\tau(t)) - \Lambda(\tau(s)) = t-s.
}
%It follows that $M_t$ is a simple counting process with independent and stationary increments and unit rate.  
\end{proof}
\begin{cor}
Let $\Lambda(t)$ be a continuous non-decreasing function. 
Then, $S_1, S_2, \dots$ are the arrival instants in a non-stationary Poisson process $N_t$ with mean function $\Lambda(t) = \E N_t$ iff $\Lambda(S_1), \Lambda(S_2), \dots$ are the arrivals instants of a stationary Poisson process of unit rate.
\end{cor}
\begin{proof}
We can write the $n$th arrival instant  $S'_n$ of unit-rate stationary Poisson process $M_t$, in terms of the $n$th arrival instant $S_n$ of non-stationary Poisson process $N_t$ as
\begin{align*}
S'_n &= \inf\set{t > 0: {\tau(t)} \ge S_n} = \sup\set{ t > 0: \Lambda(S_n) \ge t} = \Lambda(S_n).
\end{align*}
\end{proof}
This corollary implies that $S_n \in (s,t]$ if and only if $\Lambda(S_n) \in (\Lambda(s), \Lambda(t)]$. 
Therefore, the number of arrivals in $(s,t]$ equals the number of arrivals for unit-rate stationary Poisson process in $(\Lambda(s), \Lambda(t)]$. 
Hence, we conclude that for $\Lambda(s,t] = \Lambda(t) - \Lambda(s)$
\begin{align*}
P\set{N_{t}-N_s = k} &= e^{-\Lambda(s,t]}\frac{\Lambda(s,t]^k}{k!},~k \in \Z_+.
\end{align*}
We will see that the inter-arrival times $T_n \triangleq S_n  - S_{n-1}$ for the non-stationary Poisson process $N_t$ 
%, defined as 
%\meq{2}{
%&T_0 = 0, &&T_n = S_n - S_{n-1},~n \in \N,
%}
are not independent anymore. 
\begin{prop} 
For a non-stationary Poisson process with continuous mean function $\Lambda(t)$, we have
\begin{align*}
\E[\indicator{T_{n+1} > t} \given S_1, S_2, \dots, S_n]&= \exp\left(-\Lambda(S_n+t) + \Lambda(S_n)\right).
\end{align*}
\end{prop}
\begin{proof}
We define events $A = \set{\Lambda(S_{n+1}) > \Lambda(S_n+t)}$ and $B = \set{\Lambda(S_{n+1}) \ge \Lambda(S_n+t)}$. 
Then, we have 
%\begin{align*}
$A \subseteq \set{T_{n+1} > t} \subseteq B$.
%\end{align*}
Hence, we can write 
\EQ{
\E[\indicator{A} \given S_1, S_2, \dots, S_n] \le \E[\indicator{T_{n+1} > t }\given S_1, S_2,\dots, S_n].
}
The arrival instants $S_1, \dots, S_n$ determine $\Lambda(S_1), \dots, \Lambda(S_n), \Lambda(S_n+t)$. 
Further, since $\Lambda(S_{n+1}) - \Lambda(S_n)$ is the inter-arrival time of the stationary Poisson process $M_t$, 
it is independent of $S_1,S_2,\dots, S_n$
\EQ{
\E[\indicator{A}\given S_1,\dots, S_n] = \E[\indicator{\Lambda(S_{n+1})-\Lambda(S_n) > \Lambda(S_n+t) - \Lambda(S_n)} \given S_1,\dots,S_n] = \exp\left(-\Lambda(S_{n}+t)+\Lambda(S_n)\right).
}
Result follows from the continuity of the exponential distribution. 
\end{proof}

\end{document}